\documentclass[11pt,a4paper]{article}

% ── Packages (only those available on this installation) ──────────────────────
\usepackage[margin=2.5cm]{geometry}
\usepackage{amsmath,amssymb}
\usepackage{graphicx}
\usepackage{booktabs}
\usepackage{longtable}
\usepackage{array}
\usepackage{xcolor}
\usepackage{hyperref}
\usepackage{caption}
\usepackage{subcaption}
\usepackage{enumitem}
\usepackage{microtype}
\usepackage{listings}
\usepackage{float}
\usepackage{placeins}

% Text extraction: map every glyph, ligatures included, to Unicode (review of 2026-09-28)
\input{glyphtounicode}
\pdfgentounicode=1

% ── Figure path ───────────────────────────────────────────────────────────────
\graphicspath{{figures/}}

% ── Hyperref setup ────────────────────────────────────────────────────────────
\hypersetup{
  colorlinks = true,
  linkcolor  = blue!70!black,
  citecolor  = green!50!black,
  urlcolor   = blue!70!black,
}

% ── Custom commands ───────────────────────────────────────────────────────────
\newcommand{\numu}{$\nu_\mu$}
\newcommand{\numubar}{$\bar{\nu}_\mu$}
\newcommand{\ccpip}{$\mathrm{CC}\pi^{\pm}$}
\newcommand{\pmu}{$p_\mu$}
\newcommand{\ppi}{$p_\pi$}
\newcommand{\costhmu}{$\cos\theta_\mu$}
\newcommand{\costhpi}{$\cos\theta_\pi$}
\newcommand{\thmupi}{$\theta_{\mu\pi}$}
\newcommand{\uboone}{MicroBooNE}
\newcommand{\GeVc}{\,\text{GeV}\!/c}
\newcommand{\pot}{\,\text{POT}}
\newcommand{\cm}{\,\text{cm}}

% inline math shorthand so they work inside math mode too
\newcommand{\pmuM}{p_\mu}
\newcommand{\ppiM}{p_\pi}
\newcommand{\cthmuM}{\cos\theta_\mu}
\newcommand{\cthpiM}{\cos\theta_\pi}
\newcommand{\thmupM}{\theta_{\mu\pi}}

% ── Listings (C++ snippets) ───────────────────────────────────────────────────
\lstset{
  language=C++,
  basicstyle=\ttfamily\small,
  keywordstyle=\bfseries\color{blue!70!black},
  commentstyle=\itshape\color{gray},
  backgroundcolor=\color{gray!10},
  frame=single,
  framerule=0.4pt,
  numbers=none,
  breaklines=true,
}

% ── Table column helpers ──────────────────────────────────────────────────────
\newcolumntype{C}[1]{>{\centering\arraybackslash}p{#1}}
\newcolumntype{R}[1]{>{\raggedleft\arraybackslash}p{#1}}
\newcolumntype{L}[1]{>{\raggedright\arraybackslash}p{#1}}

% Let TeX stretch spacing on hard-to-break lines (long \texttt paths/code) to keep
% text out of the margin.
\emergencystretch=3em

% ─────────────────────────────────────────────────────────────────────────────


% Cross-document references: this document collects the approval request and the prerequisites that were
% moved out of the analysis note, the proton-tagged note and the technical supplement on 2026-09-28.
\usepackage{xr}
\externaldocument[N-]{analysis_note}
\externaldocument[P-]{proton_tagged_note}
\externaldocument[S-]{technical_supplement}
\newcommand{\noteref}[2]{#1~\ref{N-#2} of the analysis note}
\newcommand{\ptref}[2]{#1~\ref{P-#2} of the proton-tagged note}
\newcommand{\suppref}[2]{#1~\ref{S-#2} of the technical supplement}
\newcommand{\inclref}[2]{\noteref{#1}{#2}}
% status_vocab.tex -- the four result-status labels, shared by the analysis note, the proton-tagged note and
% the technical supplement (decision of 2026-09-27). Use only these; report/tools/check_documents.py refuses
% the older overlapping labels ("proposed secondary", "secondary scope", ...). 2026-09-28: "Secondary result proposed
% for approval" shortened to "Secondary result" (the analysis documents describe the analysis as it stands).
\providecommand{\statPrimary}{Primary result}
\providecommand{\statSecondary}{Secondary result}
\providecommand{\statSecondaryPl}{Secondary results}
\providecommand{\statValidation}{Validation-only product}
\providecommand{\statNotReported}{Not reported}


\begin{document}
\begin{center}
{\LARGE\bfseries Approval request and prerequisites\par}
\vspace{0.4cm}
{\large Measurement of $\nu_\mu+\bar\nu_\mu$ charged-current single-pion production on argon in the NuMI
beam with the \uboone{} detector\par}
\vspace{0.4cm}
{\large Ishanee Pophale and Jarek Nowak\par}
\vspace{0.2cm}
{\normalsize\itshape Lancaster University\par}
\end{center}
\vspace{0.4cm}

This document contains the decision requested from the collaboration, the status of every result with
the prerequisites that apply to it, the validation ladder, and the request to compare the control
regions with beam-on data. The analysis is described in the analysis note, the proton-tagged note and
the technical supplement.

\tableofcontents

% =============================================================================
\section{Decision requested}
\label{sec:ar_decision}
% =============================================================================

% decision_box.tex -- the decision requested from the collaboration. Included verbatim at the beginning of
% the analysis note and of the proton-tagged note (and in the technical supplement), so the request cannot
% differ between documents. The prerequisite identifiers refer to Table~\ref{tab:prerequisites}, which each
% document includes from shared/prerequisites.tex.
\begin{center}
\fbox{\parbox{0.95\textwidth}{\small
\textbf{Decision requested from the collaboration}
\begin{enumerate}\itemsep1pt
  \item Approve the inclusive CC$\pi^\pm$ signal definition, selection, extraction and uncertainty model,
    and the reported inclusive observables, each as a \statPrimary{}: $d\sigma/dp_\mu$,
    $d\sigma/d\cos\theta_\mu$, $d\sigma/d\cos\theta_\pi$, $d\sigma/d\theta_{\mu\pi}$, the cross section in
    three $p_\pi$ regions and the one-bin total cross section, each in the FHC, RHC and combined
    configurations.
  \item Do not yet authorise opening of the inclusive signal region: the prerequisites U1--U6 of
    Table~\ref{tab:prerequisites} remain to be completed.
  \item Approve the following, each as a \statSecondary{}: the proton-tagged one-bin total cross
    section and the
    proton-tagged $d\sigma/d\theta_p$ and $d\sigma/d\theta_{\pi p}$ (FHC, RHC, combined), and the
    two-dimensional results $d^2\sigma/d\theta_p\,d\delta p_T$ (proton-tagged sample) and
    $d^2\sigma/d\cos\theta_\pi\,d\cos\theta_\mu$ (inclusive sample) in the combined configuration only,
    subject to the prerequisites listed for each in Table~\ref{tab:prerequisites}.
  \item Treat every other proton-tagged extraction ($p_p$, $W_\mathrm{had}$, $\delta p_T$,
    $\delta\alpha_T$, $\delta\phi_T$, $p_n$ and the muon and pion observables on the proton-tagged
    sample) as a \statValidation{}.
  \item Do not report the differential $W_{\pi p}$ result, the other two-dimensional pairs, or the FHC and
    RHC extractions of the two two-dimensional results (\statNotReported{}).
\end{enumerate}}}
\end{center}


% =============================================================================
\section{Status of the results}
\label{sec:ar_status}
% =============================================================================

Tables~\ref{tab:inventory_prereq} and~\ref{tab:pt_status_prereq} give the status of every inclusive and
proton-tagged product and the prerequisites of Table~\ref{tab:prerequisites} that apply to it.

% generated by report/tools/inventory.py from report/notes/status.tsv -- do not edit
{\footnotesize\setlength{\tabcolsep}{4pt}
\begin{longtable}{@{}p{6.1cm}ccc p{3.3cm} p{2.2cm}@{}}
\caption{Status of the inclusive products. ``yes'' marks the configurations in which a product is reported. The prerequisites are those of Table~\ref{tab:prerequisites}.}
\label{tab:inventory_prereq}\\
\toprule
Product & FHC & RHC & Comb. & Status & Prerequisites \\
\midrule
\endfirsthead
\toprule
Product & FHC & RHC & Comb. & Status & Prerequisites \\
\midrule
\endhead
\bottomrule
\endfoot
$d\sigma/dp_\mu$, $d\sigma/d\cos\theta_\mu$, $d\sigma/d\cos\theta_\pi$, $d\sigma/d\theta_{\mu\pi}$, $\sigma$ in three $p_\pi$ regions & yes & yes & yes & \statPrimary & U1--U7 \\
One-bin total cross section & yes & yes & yes & \statPrimary & U1--U3, U5--U7 \\
$d^2\sigma/d\cos\theta_\pi\,d\cos\theta_\mu$ & no & no & yes & \statSecondary & U1--U3, U5--U7, R1 \\
$d\sigma/dp_\pi$ in five bins &  &  &  & \statNotReported &  \\
$d\sigma/d\theta_\mu$ &  &  &  & \statNotReported &  \\
Two-dimensional $\cos\theta_\pi\times\theta_{\mu\pi}$, $\cos\theta_\mu\times p_\mu$, $\cos\theta_\pi\times p_\pi$; FHC and RHC $\cos\theta_\pi\times\cos\theta_\mu$ &  &  &  & \statNotReported &  \\
\end{longtable}}


% generated by report/tools/inventory.py from report/status.tsv -- do not edit
{\footnotesize\setlength{\tabcolsep}{4pt}
\begin{longtable}{@{}p{6.1cm}ccc p{3.3cm} p{2.2cm}@{}}
\caption{Status of the proton-tagged products. ``yes'' marks the configurations in which a product is reported. The prerequisites are those of Table~\ref{tab:prerequisites}.}
\label{tab:pt_status_prereq}\\
\toprule
Product & FHC & RHC & Comb. & Status & Prerequisites \\
\midrule
\endfirsthead
\toprule
Product & FHC & RHC & Comb. & Status & Prerequisites \\
\midrule
\endhead
\bottomrule
\endfoot
One-bin total cross section & yes & yes & yes & \statSecondary & U1--U3, U5--U7, P1 \\
$d\sigma/d\theta_p$, $d\sigma/d\theta_{\pi p}$ & yes & yes & yes & \statSecondary & U1--U3, U5--U7, P1, P2 \\
$d^2\sigma/d\theta_p\,d\delta p_T$ & no & no & yes & \statSecondary & U1--U7, P1--P3, R1 \\
$d\sigma/dp_p$ & yes & yes & yes & \statValidation & P1 \\
$d\sigma/dW_\mathrm{had}$ (two regions) & yes & yes & yes & \statValidation & P1, P4 \\
$d\sigma/d\delta p_T$, $d\sigma/dp_n$ & yes & yes & yes & \statValidation & P1, P3, P4 \\
$d\sigma/d\delta\alpha_T$, $d\sigma/d\delta\phi_T$ & yes & yes & yes & \statValidation & P1 \\
$p_\mu$, $p_\pi$ (two regions), $\cos\theta_\mu$, $\cos\theta_\pi$, $\theta_{\mu\pi}$ on the proton-tagged sample & yes & yes & yes & \statValidation & P1 \\
$d\sigma/dW_{\pi p}$ &  &  &  & \statNotReported &  \\
Two-dimensional $\theta_p\times W_{\pi p}$, $\theta_{\pi p}\times W_{\pi p}$, $\theta_{\pi p}\times\delta p_T$; FHC and RHC $\theta_p\times\delta p_T$ &  &  &  & \statNotReported &  \\
\end{longtable}}


% =============================================================================
\section{Validation ladder}
\label{sec:ar_ladder}
% =============================================================================

Table~\ref{tab:validation_ladder} lists the validation steps with their status and criteria. Closure, the
ensembles and the reweighted-model test are complete for the primary and the secondary results; the
independent-generator closure (U3), the control-region comparison (U2, and P1 for the proton
mis-assignment) and the prerequisites for opening the signal region are not.

% validation_ladder.tex -- the validation steps of the extraction, their status, criterion, worst observed
% case and blocking level; shared verbatim by the analysis note, the proton-tagged note and the technical
% supplement (2026-09-27). Prerequisite identifiers refer to Table~\ref{tab:prerequisites}.
{\footnotesize\setlength{\tabcolsep}{4pt}
\begin{longtable}{@{}p{1.9cm}p{2.7cm}p{2.6cm}p{2.7cm}p{3.1cm}p{1.4cm}@{}}
\caption{Validation ladder. Each step asks a question that the steps before it cannot answer; closure under
  the nominal model is necessary but not sufficient, since a result can close while being determined largely
  by the modelled migrations or by the regularisation, which the migration criterion and the trace of $A_C$
  address separately (\inclref{Sec.}{sec:binning_rule}).}
\label{tab:validation_ladder}\\
\toprule
Step & Question & Status & Criterion & Worst case observed & Blocks \\
\midrule
\endfirsthead
\toprule
Step & Question & Status & Criterion & Worst case observed & Blocks \\
\midrule
\endhead
\bottomrule
\endfoot
1. Implementation closure & Does the chain return the truth of pseudo-data thrown from the model of the
  response? & Done for every extraction &
  Unfolded result consistent with the $A_C$-smeared truth; integral ratio within the statistical spread &
  Inclusive: $\chi^2\le2.39/8$, integral ratios $0.946$--$1.047$. Proton-tagged: $\chi^2\le6.01/4$ (RHC
  $\cos\theta_\pi$), ratios $0.93$--$1.18$; two-dimensional: cell ratios $0.70$--$1.36$ & Unblinding \\
2. Statistical calibration & Are the quoted statistical uncertainties correctly sized? &
  Done: six inclusive and three proton-tagged extractions, $100$ pseudo-experiments each. Outstanding for
  $\theta_p$, $\theta_{\pi p}$ and the two-dimensional results (V1) &
  Pull width consistent with one within the ensemble precision ($\pm13\%$ at $100$ members); coverage
  $68\%$ and $95\%$ &
  Inclusive: widths $0.95$--$1.07$, $68\%$ coverage $63$--$73\%$. Proton-tagged: widths $0.94$--$1.08$,
  $62$--$72\%$ & Unblinding \\
3. Parametric robustness & Does the extraction follow a different model that reweighting can reach? &
  Done: GENIE multisim universe 545 and $\Delta\to N\pi$ angular, inclusive FHC and RHC, proton-tagged FHC.
  Outstanding for $\theta_p$, $\theta_{\pi p}$ and the two-dimensional results (V1) &
  Per-bin bias against the per-bin total uncertainty; a bias above $1\sigma$ is flagged and treated &
  Inclusive: $1.18\sigma$ in one sparse bin (FHC $\cos\theta_\mu$). Proton-tagged: $-1.14\sigma$ and
  $-1.16\sigma$ in the low $\delta p_T$ and $p_n$ bins (P3) & Unblinding \\
4. Structural robustness & Does it close on pseudo-data from a generator that reweighting cannot reach? &
  Outstanding (U3) & Criterion of U3 & n/a & Unblinding \\
5. Background validation & Do control-region data support the background model? &
  Outstanding (U2; P1 for the proton mis-assignment) &
  Frozen procedure of the control regions & n/a & Unblinding \\
6. Signal-region opening & Are the normalisation and covariance prerequisites complete? &
  Outstanding & U1--U6 complete & n/a & n/a \\
\end{longtable}}


% =============================================================================
\section{Prerequisites}
\label{sec:ar_prereq}
% =============================================================================

\paragraph{Inclusive results.} The inclusive signal region can be opened once U1--U6 are complete; U7 is
required before publication, and U4 applies to the results that use the muon momentum. The flux
confirmation (U1) concerns the three integrated fluxes of \noteref{Table}{tab:flux}: the NuMI flux group
is asked to confirm in writing that they are the PPFX-corrected central values that correspond to the
per-event PPFX weights of the simulation (\noteref{Sec.}{sec:flux_status}). The beam-on exposure check
(U5) replaces the recorded exposures of \noteref{Table}{tab:pot} by values summed over the runs contained
in each beam-on file. The size of the MCS momentum-scale variation, $\pm5\%$, is to be confirmed by a
comparison of the MCS and range momenta of contained muons in data (U4; \noteref{Sec.}{sec:mcs}). The
clean-cache reproduction of \suppref{Sec.}{sec:pipeline} is to be repeated on the release to be unblinded
(U6), and the three terms outside the official covariance (\noteref{Sec.}{sec:notincov}) are to be
included before publication (U7).

\paragraph{Independent-generator closure (U3).} The test uses a NuWro overlay with the full detector
simulation for FHC and RHC. The overlay is first validated against standalone NuWro on the same flux: its
signal rate per POT must agree with the standalone prediction. Pseudo-data thrown from it are then
unfolded with the nominal response for every primary and secondary result and compared with the
$A_C$-smeared NuWro truth. The test passes if every bin holding at least $10\%$ of the largest truth lies
within $1\sigma$ of its total uncertainty and every integral within its uncertainty; the signal region is
not opened unless it passes. The one existing NuWro overlay, for Run-1 FHC, cannot serve
(\noteref{Sec.}{sec:altmodel}).

\paragraph{Proton-tagged results.} The proton-tagged secondary results require the inclusive
prerequisites U1--U3 and U5--U7, U4 for $\theta_p\times\delta p_T$, and the proton-tagged items listed
for each in Table~\ref{tab:pt_status_prereq}. The inclusive control regions do not validate the dominant
proton-candidate misidentification of this sample (\ptref{Sec.}{sec:pt_limitations}). Until an
inverted-proton-PID region is defined and validated (P1), no proton-tagged result is unblinded, and
$\theta_p$, $\theta_{\pi p}$ and $\theta_p\times\delta p_T$, whose migrations are dominated by the
proton-candidate assignment, carry the additional term of P2. The reconstruction of $W_\mathrm{had}$ and
$p_n$ needs a truth-level validation for pion production or a redefinition of the observables as the
proxies they are (P4). The model sensitivity of the low $\delta p_T$ and $p_n$ bins needs an explicit
model term in the covariance, the removal of those bins, or an independent-generator test that shows no
such shift (P3).

\paragraph{Reporting decision (R1).} Every one-dimensional result is quoted per beam mode and for the
combination. The responses of the two two-dimensional results can be inverted only in the combined
configuration, and whether a result available only there is acceptable is a collaboration decision.

% prerequisites.tex -- the numbered prerequisites of the analysis, shared verbatim by the analysis note, the
% proton-tagged note and the technical supplement (decision of 2026-09-27). The identifiers are used by
% report/status.tsv, the status tables and the decision box; report/tools/check_documents.py checks that every
% identifier used anywhere is defined here. Blocking levels: unblinding (the signal region of the results
% concerned is not opened before the item is complete), publication, approval (the result is not approved),
% promotion (a validation-only product cannot become a proposed result).
{\footnotesize\setlength{\tabcolsep}{4pt}
\setlength{\LTcapwidth}{\textwidth}
\begin{longtable}{@{}p{0.6cm}p{2.5cm}p{5.9cm}p{2.8cm}p{2.4cm}@{}}
\caption{Prerequisites, their criteria and what each blocks. U: all results; P: proton-tagged results;
  V: results proposed since the last validation round; R: reporting decisions.}
\label{tab:prerequisites}\\
\toprule
ID & Prerequisite & Criterion & Owner; evidence & Blocks \\
\midrule
\endfirsthead
\toprule
ID & Prerequisite & Criterion & Owner; evidence & Blocks \\
\midrule
\endhead
\bottomrule
\endfoot
U1 & Absolute flux normalisation &
  Written confirmation that the integrated $\nu_\mu+\bar\nu_\mu$ fluxes above $60$~MeV used for the
  normalisation (FHC $6.81159$, RHC $6.44646$, combined $6.59691\times10^{-10}$~cm$^{-2}$~POT$^{-1}$) are the
  PPFX-corrected central values that correspond to the per-event PPFX weights &
  NuMI flux group (flux contact); written statement & Unblinding \\
U2 & Control regions &
  The four inclusive control regions compared with beam-on data under the frozen procedure, and any
  discrepancy treated as that procedure prescribes &
  Analysers and conveners; control-region comparison & Unblinding \\
U3 & Independent-generator closure &
  A NuWro overlay with the full detector simulation, validated against standalone NuWro on the same flux,
  unfolded with the nominal response: every bin holding at least $10\%$ of the largest truth within
  $1\sigma$ of its total uncertainty, and every integral within its uncertainty &
  Analysers, with the production group for the overlay; closure tables for every primary and proposed
  result & Unblinding \\
U4 & MCS momentum scale &
  Comparison of the MCS and range momenta of contained muons in data and simulation: the difference of
  the MCS scale within the $\pm5\%$ of the covariance term, or the term enlarged to the observed
  difference &
  Analysers; MCS-to-range comparison & Unblinding of results that use the muon momentum \\
U5 & Beam-on exposure &
  POT and trigger counts of every beam-on file summed over the runs the file contains, and checked against
  the beam database to better than $1\%$ &
  Analysers, with the beam-data contact; per-run exposure table & Unblinding \\
U6 & Reproduction &
  The result set rebuilt from empty caches reproduces every released bin to $10^{-9}$ &
  Analysers; rebuild log & Unblinding \\
U7 & Terms outside the covariance &
  The beamline-geometry flux term ($1.6$--$3.3\%$), the beam-off gate-ratio term and the Run-2 beam-off
  stand-in term included in the released covariance &
  Analysers; released covariance & Publication \\
\midrule
P1 & Proton mis-assignment control region &
  An inverted-proton-PID control region at exactly one pion candidate, validated with pseudo-data and
  compared with beam-on data under the frozen procedure; the rate of wrong proton candidates in the
  simulation consistent with the data within its uncertainty, or corrected under that procedure &
  Analysers; region definition and comparison & Unblinding of every proton-tagged result \\
P2 & Wrong-candidate term for the proton angles &
  The fraction of wrong proton candidates varied by its uncertainty from P1, and the resulting change of
  every bin of $\theta_p$, $\theta_{\pi p}$ and $\theta_p\times\delta p_T$ added to their covariance &
  Analysers; covariance term & Unblinding of $\theta_p$, $\theta_{\pi p}$, $\theta_p\times\delta p_T$ \\
P3 & Model sensitivity of the low-imbalance bins &
  With an explicit model term in the covariance, the alternative-model bias of the low $\delta p_T$ and
  $p_n$ bins below $1\sigma$; or those bins removed; or U3 showing no such shift &
  Analysers; covariance term or closure & Unblinding of $\theta_p\times\delta p_T$; promotion of
  $\delta p_T$, $p_n$ \\
P4 & Reconstruction validity of $W_\mathrm{had}$ and $p_n$ &
  For pion production, a truth-level study showing that each reconstructed proxy follows the quantity it
  is named for (hadronic invariant mass, initial-nucleon momentum) with a bias smaller than the bin width
  in every bin; otherwise the observable redefined as the proxy, with the same prescription at truth
  level &
  Analysers; truth-level study & Promotion of $W_\mathrm{had}$, $p_n$ \\
\midrule
V1 & Validation of the newly proposed results &
  Ensembles (pull widths consistent with one within the ensemble precision, nominal coverage) and the
  alternative-model test (per-bin bias below $1\sigma$, or flagged and treated) for $\theta_p$,
  $\theta_{\pi p}$ and the two two-dimensional results &
  Analysers; ensemble and closure tables & Unblinding of those results \\
R1 & Results in the combined configuration only &
  A collaboration decision on quoting the two two-dimensional results only in the combined
  configuration, since their FHC and RHC responses cannot be inverted &
  Conveners; decision recorded & Approval of the two-dimensional results \\
\end{longtable}}


% =============================================================================
\section{Control regions: request and sign-off}
\label{sec:ar_cr}
% =============================================================================

The decision requested of the collaboration is permission to inspect beam-on data in four
reconstruction-level control regions of the inclusive $\mathrm{CC}\,1\mu1\pi Xp$ selection while the
signal region stays blind. \noteref{Sec.}{sec:sidebands} gives the region definitions and
\suppref{Sec.}{sec:request} their roles.

\paragraph{Analysis state.} Repository tag \texttt{cr-freeze-2026-09-06};
pseudo-data: one Poisson throw per run of the central-value prediction with the unified
weight rule (\noteref{Sec.}{sec:mc}); binning: the released edges of \noteref{Appendix}{app:binning}
(five-bin forward-split $\cos\theta_\pi$, the $p_\pi$ binning then in use); normalisation: per-run
exposure, per-run beam-off gate ratios, per-configuration flux; detector variations:
native Run-4 FHC and RHC samples; data release: \texttt{report/data\_release/}, checksum
manifest \texttt{MANIFEST.sha256} (\texttt{b8ebc69e3bc91434}).

\paragraph{Request.}
We request approval to inspect beam-on data in four control regions: CC$0\pi$,
multi-$\pi$, $\pi^0$ and large muon--pion opening angle (called ``cosmic'' below for
brevity: it is cosmic-enriched, $34$--$41\%$ beam-off, not cosmic-pure).
Throughout, ``beam-off'', ``EXT'' and ``cosmic data'' name the same data-driven
component, the beam-off trigger sample scaled by gate counts; ``signal contamination''
means the inclusive signal unless the proton-tagged selection is named. Each region is
anchored by the inversion of one nominal signal requirement; requirements that become
undefined once that inversion is made, or that would sculpt the targeted control sample,
are dropped as listed in \suppref{Table}{tab:cr_defs}. Each region is disjoint from the
signal region at reconstruction level (demonstrated event by event,
\suppref{Sec.}{sec:cr_overlap}); their expected signal contaminations are $8$--$24\%$. The
signal region remains blind. Before the regions are opened, the selections, the
run-dependent normalisation, the beam-off scaling, the displayed observables and binning,
the covariance construction, the goodness-of-fit tests and the permitted responses to a
discrepancy are frozen as written in \suppref{Sec.}{sec:cr_protocol}. Control-region data will be
compared with the complete prediction (neutrino MC $+$ beam-off $+$ dirt) separately in
FHC and RHC, with the combined configuration as a cross-check. Any discrepancy beyond
the pre-declared criteria requires a documented treatment and renewed approval before
the signal region is opened. The proton-tagged control regions are not part of
this request (\suppref{Sec.}{sec:cr_1p_scope}).

\paragraph{Points the collaboration is asked to ratify.}
(i)~The $\pi^0$ definition without the final multiplicity cut (\suppref{Table}{tab:cr_defs});
the comparison with the version retaining it is in \suppref{Sec.}{sec:cr_pi0_final} and will
not be used to modify the region after beam-on inspection. (ii)~The $\pi^0$
target class CC$+$NC~$\pi^0$, adopted by the analysers (the NC$\pi^0$
events the shower veto is also meant to remove had been counted under ``NC and
$\nu_e$CC''; $\nu_\mu$CC~$\pi^0$ alone is $46\%$ of the region), together with the
region's diagnostic-only status for this opening. Purity is defined throughout as the
target class's share of the neutrino-MC prediction in the region, beam-off and
dirt being separately constrained components; both the class definition and the
denominator were settled after the purities were known, which is one reason the region
is diagnostic-only. (iii)~The roles of the other three regions as stated above, and the
frozen procedure of \suppref{Sec.}{sec:cr_protocol}. The collaboration is also asked to ratify the frozen $\pi^0$ definition
with the comparison of \suppref{Sec.}{sec:cr_pi0_final} in view.

\paragraph{Approval within the procedure.} The analysers propose the procedure of
\suppref{Sec.}{sec:cr_protocol} for ratification with the approval. It is frozen, and any departure from
it is itself a change requiring approval. The exclusion of a run period (item~4(i)) requires a new frozen
release and explicit approval. The $\pi^0$ region is diagnostic only for this opening: approval to
inspect it does not authorise a fitted correction, and a fitted $\pi^0$ correction may be proposed only
after the pre-fit result is documented, with the fit model, the leverage of the fixed components, the
transfer-factor uncertainty and its floor reviewed by the collaboration before the correction is
propagated to the signal region; a CC$0\pi$ correction is pre-authorised under the rule of item~4(ii).
If the post-treatment test still fails (item~4(iv)), the background model is revised and renewed
collaboration approval is required before the signal region is opened. Conditioning the prediction on
the control-region totals, which on an Asimov dataset reduces the uncertainty of the one-bin total from
$35.8\%$ to $11.3\%$ (\noteref{Sec.}{sec:cr_procedure}), would change the role of the control regions and
requires a collaboration decision.

\paragraph{Sign-off.} The analysis conveners, the flux and detector-systematics contacts, and the
collaboration review committee, on the basis of a short report containing the comparisons of item~1 and
the response chosen from item~4 of the procedure.

\paragraph{Proton-tagged regions.} The proton-tagged control regions are not part of this request: they
apply the proton tag rather than inverting it and do not constrain the proton mis-tag background
(\suppref{Sec.}{sec:cr_1p_scope}), so the proton-tagged control-region programme is not ready for
equivalent approval.

\paragraph{Approval checklist.}
\begin{center}\small
\begin{tabular}{p{7.2cm}p{4.6cm}p{2.6cm}}
\toprule
Item & Status / evidence & Sign-off \\
\midrule
\multicolumn{3}{l}{\emph{Required before the control regions are opened}}\\
\midrule
Control-region definitions frozen & \suppref{Table}{tab:cr_defs} & \\
Signal-region overlap zero, event by event & \suppref{Sec.}{sec:cr_overlap} & \\
Plot list and binning frozen & item~1 of \suppref{Sec.}{sec:cr_protocol} & \\
Pre-fit covariance, global statistic and thresholds frozen & items~2--3 of \suppref{Sec.}{sec:cr_protocol}; null calibration \suppref{Fig.}{fig:cr_null_calibration} & \\
Response to a discrepancy frozen & item~4 of \suppref{Sec.}{sec:cr_protocol} & \\
$\pi^0$ definition and target class ratified; diagnostic-only status recorded & points (i)--(ii), this section & conveners \\
Flux and detector-systematics contacts have reviewed their parts & pending & contacts \\
Repository tag and release checksum recorded & analysis state, this section & \\
\midrule
\multicolumn{3}{l}{\emph{Required later, before the signal region is opened}}\\
\midrule
Control-region results reviewed; any discrepancy treated under \suppref{Sec.}{sec:cr_protocol} & pending & conveners \\
Model-distortion pseudo-data closure (reweighted GENIE, FHC) & done, \noteref{Sec.}{sec:validation} & \\
Model-distortion pseudo-data closure, proton-tagged family (FHC) & done, \ptref{Sec.}{sec:pt_altmodel}; low-imbalance $\delta p_T$/$p_n$ bins biased $1.1$--$1.2\sigma$ under GENIE u545 & \\
Statistical-coverage ensembles, proton-tagged family ($W_{\rm had}$, $\delta p_T$, $p_n$) & done, \ptref{Sec.}{sec:pt_ensemble}; pull widths $0.94$--$1.08$, integral offsets within $0.7\%$ & \\
Ensembles and alternative-model closure of the results proposed since then, $\theta_p$, $\theta_{\pi p}$, $\theta_p\times\delta p_T$ and $\cos\theta_\pi\times\cos\theta_\mu$ (prerequisite V1) & done, \ptref{Sec.}{sec:pt_ensemble}, \ptref{Sec.}{sec:pt_altmodel}, \noteref{Sec.}{sec:incl_2d}; pull widths $0.94$--$1.04$, largest bias $0.79\sigma$ & \\
Independent-generator pseudo-data closure (prerequisite U3): a NuWro overlay with the full detector simulation, validated against standalone NuWro on the same flux, unfolded with the nominal response & pending: the existing NuWro overlay failed validation against the standalone prediction and is not used, \noteref{Sec.}{sec:limitations} & \\
Flux constants confirmed as PPFX-corrected central values & pending, \noteref{Sec.}{sec:flux_status} & flux contact \\
MCS muon-momentum scale ($\pm5\%$, data side) evaluated on $p_\mu$ in FHC, RHC and combined and released as \texttt{cov\_MCSscale.txt} & done, \noteref{Sec.}{sec:mcs}; at most $0.56\sigma$ (FHC open top bin), $0.42\sigma$ (combined) and $0.30\sigma$ (RHC); exactly zero for the other inclusive observables, which carry no reco-$p_\mu$ cut; computed outside the framework and part of the official covariance (\texttt{cov\_total\_plusMCS.txt}) and of every summary table; size to be confirmed by prerequisite U4; proton-tagged observables evaluated (at most $0.62\sigma$, FHC $p_\mu$; \ptref{Sec.}{sec:pt_mcs}) & \\
Final empty-cache reproduction of the release to be unblinded & pending, \suppref{Sec.}{sec:pipeline} & \\
Inverted pion-PID region for the CC$0\pi$ misidentification transfer & pending & \\
Inverted-proton-PID region before any proton-tagged result & pending, \suppref{Sec.}{sec:cr_1p_scope} & \\
\bottomrule
\end{tabular}
\end{center}

\end{document}
